\section{Interpretabilidad y Explicabilidad de Inteligencia Artificial}
Hasta el momento se ha construido una base para entender el funcionamiento de redes neuronales y 
redes convolucionales desde un punto de vista técnico. Sin embargo, la creciente adopción de sistemas 
de inteligencia artificial en distintas industrias y en la vida cotidiana ha generado la necesidad 
de desarrollar modelos más transparentes, cuyos procesos de toma de decisiones resulten comprensibles 
para los usuarios humanos. A partir de esta necesidad surgen los conceptos de inteligencia artificial 
interpretable e inteligencia artificial explicable (\textit{Explainable Artificial Intelligence}, 
o XAI).
\medskip

La interpretabilidad de la inteligencia artificial se refiere a la capacidad de un modelo para 
proporcionar un entendimiento claro de sus mecanismos internos y predicciones. Generalmente, 
se consideran interpretables aquellos modelos cuyas arquitecturas y mecanismos son lo suficientemente 
simples como para ser comprendidos sin herramientas adicionales. A estos modelos también se les 
conoce como modelos de caja blanca o modelos transparentes. Un ejemplo de este tipo es el 
perceptrón, un modelo de red neuronal compuesto únicamente por una capa de entrada y una capa de 
salida, cuya simplicidad permite describir matemáticamente la relación entre las entradas y 
las salidas. Otros ejemplos de modelos transparentes incluyen la regresión lineal, 
la regresión logística y los árboles de decisiones.
\medskip

Sin embargo, como se describió en secciones anteriores, las redes neuronales modernas poseen 
arquitecturas considerablemente más complejas: una sola capa puede contener miles de parámetros. 
Describir matemáticamente la relación entre las entradas y salidas de estos modelos puede implicar 
la consideración de miles, o incluso millones, de parámetros. Debido a ello, el entendimiento de 
estos sistemas suele abstraerse a niveles más altos, mientras que su funcionamiento interno detallado 
permanece difícil de interpretar. Por esta razón, modelos de aprendizaje profundo como las redes 
convolucionales y los grandes modelos de lenguaje (\textit{Large Language Models}, o LLMs) suelen 
clasificarse como modelos de caja negra o post-hoc \parencite{dwivedi2023explainable}.
\medskip

Como respuesta a la necesidad de comprender mejor los procesos de toma de decisiones de los sistemas 
de inteligencia artificial, surge el campo de la inteligencia artificial explicable 
(XAI), cuyo propósito es desarrollar métodos que permitan mejorar la interpretabilidad y 
comprensibilidad de los modelos, ya sea mediante el diseño de arquitecturas inherentemente 
transparentes o a través de herramientas externas orientadas a explicar su comportamiento y 
predicciones \parencite{arrieta2020explainable}.

\subsection{Técnicas Agnósticas y Específicas al Modelo}
La explicabilidad post-hoc tiene como objetivo los modelos que no son inmediatamente interpretables
(los modelos de caja negra o post-hoc), por lo que requieren de técnicas adicionales para mejorar su
interpretabilidad. Estas técnicas se pueden clasificar en técnicas especificas y agnósticas al modelo.
\medskip

\subsubsection{Técnicas Agnósticas al Modelo}
Las técnicas agnósticas al modelo están diseñadas para explicar el proceso de predicción de 
distintos tipos de modelos independientemente de su arquitectura interna. Estas técnicas pueden 
clasificarse, de manera general, en métodos basados en modelos sustitutos simplificados, 
métodos de estimación de relevancia de características y métodos basados en visualización \parencite{arrieta2020explainable}.
\medskip

Muchas de las técnicas basadas en simplificación del modelo se fundamentan en la extracción de 
reglas, una estrategia que busca aproximar el proceso de toma de decisiones del modelo a partir 
de pares entrada-salida. Un ejemplo de este enfoque es el algoritmo LIME 
(por sus siglas en inglés \textit{Local Interpretable Model-Agnostic Explanations}) 
\parencite{ribeiro2016should}, el cual construye un modelo lineal local que aproxima el 
comportamiento del modelo alrededor de una predicción específica. Este tipo de análisis se 
conoce como interpretabilidad local, ya que se enfoca en explicar una predicción individual. 
Por otro lado, la interpretabilidad global busca describir el comportamiento general del modelo 
completo \parencite{dwivedi2023explainable}.
\medskip

Las técnicas de estimación de relevancia de características buscan explicar las predicciones de un 
modelo mediante la asignación de niveles de importancia a las características de entrada. Un ejemplo 
de este enfoque es SHAP (por sus siglas en inglés \textit{SHapley Additive exPlanations}) 
\parencite{lundberg2017unified}, método basado en valores de Shapley que estima la contribución 
aditiva de cada característica a una predicción.
\medskip

Finalmente, las técnicas basadas en visualización son menos comunes en enfoques agnósticos al 
modelo debido a la dificultad de aplicarlas a arquitecturas generales. Sin embargo, estas técnicas 
resultan especialmente útiles en el análisis e interpretación de redes convolucionales, 
particularmente en la visualización de representaciones intermedias y activaciones internas.

\subsubsection{Técnicas Específicas al Modelo}
Estas técnicas son diseñadas para la interpretación de modelos con características y capacidades
específicas.